\documentclass[14pt]{extarticle} 

% Packages for math, layout, and styling
\usepackage[utf8]{inputenc}
\usepackage{amsmath}
\usepackage{amsfonts}
\usepackage{amssymb}
\usepackage{graphicx}
\usepackage{geometry}
\usepackage{hyperref}
\usepackage{booktabs}
\usepackage{enumitem}
\usepackage{xcolor}

\geometry{margin=1in}

% --- Document Information ---
\title{Weather AI Analytics: A Lightweight Framework for Localized Forecasting via Linear Regression and Exponentially Weighted Moving Averages}
\author{Artsiom Hancharenka, Nil Dvorakovskiy \\ \small BSU FAMCS, $3^{rd}$ group}
\date{} % Date removed as requested

\begin{document}

\maketitle

\begin{abstract}
This paper presents a computationally efficient approach to localized meteorological prediction. By shifting away from high-resource numerical models and opaque neural networks, we propose a hybrid system utilizing Simple Linear Regression (SLR) and Exponentially Weighted Moving Averages (EWMA). The system is designed to provide explainable analytics for temperature, humidity, and precipitation. Unlike deep learning models that require vast datasets, this ``Small Data'' approach utilizes a 10-day historical window to project short-term trends. Furthermore, the inclusion of an automated validation framework allows for real-time accuracy assessment, bridging the gap between complex global models and user-centric environmental monitoring.
\end{abstract}

\textbf{Keywords:} \\
Explainable AI (XAI), Persistence Modeling, Small Data, Full-Stack Analytics, Meteorological Heuristics, Linear Least Squares, Time-Series Forecasting.

\section{Motivation}
In an era of increasing climatic volatility, the demand for precise, localized weather information has never been higher. However, contemporary meteorological forecasting remains dominated by two extremes: high-level Numerical Weather Prediction (NWP) models that require massive supercomputing resources, and consumer-grade applications that offer data without transparency. There is a significant technological gap for a ``middle-ground'' solution—one that provides lightweight, computationally efficient, and explainable analytics directly to the end-user.

The motivation behind this project is to democratize the predictive process by shifting away from ``black-box'' artificial intelligence in favor of transparent statistical heuristics. By utilizing SLR and EWMA, this research proves that meaningful short-term forecasts can be generated using minimal historical datasets without specialized hardware. Furthermore, the framework empowers users to retrospectively validate model accuracy against historical ground truths, fostering scientific literacy in personal environmental monitoring.

\section{Theoretical Background}
The study delineates a bifurcated mathematical framework designed to address the distinct temporal behaviors of meteorological variables based on two scientific pillars:

\begin{enumerate}
    \item \textbf{Persistence and Momentum Theory:} In short-term meteorology, we evolve the standard persistence method by incorporating \textit{thermal momentum}. We treat the atmosphere as a fluid system with specific heat capacity and thermal inertia. If a localized environment has shown a consistent cooling trend over a week, physical energy transfer mechanisms suggest a continuation of that vector. SLR acts as the mathematical tool to quantify this momentum, allowing us to project the most likely 24-hour state.
    \item \textbf{Stochastic Volatility and Memory Decay:} While temperature changes are governed by large-scale masses, relative humidity is a high-frequency variable sensitive to immediate localized fluctuations. Humidity exhibits ``short memory,'' requiring a model that prioritizes the most recent observations. The implementation of EWMA allows the system to prioritize recent data while mathematically ``forgetting'' distant history via an exponential decay function.
\end{enumerate}

\section{System Architecture and Data Integration}
The platform is built upon a modular full-stack architecture. The system architecture ensures heavy mathematical lifting remains server-side, while the client-side remains reactive.

\subsection{Data Integration and API Orchestration}
The integrity of the forecasting engine relies on the \textbf{Open-Meteo API} ecosystem. The system translates natural language city names into precise GPS coordinates via a Geocoding layer. Once secured, the system retrieves a 10-day historical window. This training set provides the context for the SLR and EWMA algorithms. The archive API is preferred over real-time sensors to ensure a consistent, quality-controlled dataset for trend initialization.

\subsection{Backend Processing and Frontend UX}
The backend, developed with \textbf{Flask}, manages request orchestration, parses JSON payloads, and executes the mathematical logic. The \textbf{Frontend} is an asynchronous Single Page Application (SPA). It uses a card-based layout where metrics are represented through dynamic CSS3 gradients. The ``Algorithm Test'' mode creates a closed-loop system where users validate the AI's predictive logic against actual historical records.

\section{Mathematical Framework}
The predictive capability is categorized into four distinct mathematical components optimized for the physical characteristics of the target variable.

\subsection{Linear Least Squares (LLS) for Temperature Trends}
Temperature is modeled using the Least Squares Method to minimize the Residual Sum of Squares (RSS):
\begin{equation}
m = \frac{n\sum(x_iy_i) - \sum x_i \sum y_i}{n\sum x_i^2 - (\sum x_i)^2} \quad , \quad b = \frac{\sum y_i - m \sum x_i}{n}
\end{equation}
Evaluating $y = m(n) + b$ provides the projected value for the subsequent day.

\subsection{Exponentially Weighted Moving Average (EWMA) for Humidity}
To prioritize recent shifts, we implement a recursive EWMA:
\begin{equation}
S_t = \alpha \cdot Y_t + (1 - \alpha) \cdot S_{t-1}
\end{equation}
With $\alpha = 0.7$, 70\% of the prediction is based on the most recent 24 hours, ensuring the model reacts quickly to incoming fronts.

\subsection{Heuristic Thresholding and Physical Clamping}
Deterministic mathematical models often ignore physical boundaries. We implement a constraint layer:
\begin{equation}
H_{final} = \max(0, \min(100, H_{pred})) \quad , \quad P = \{0 \text{ if } P_{pred} < 0.3, \text{ else } P_{pred}\}
\end{equation}
Linear models can mathematically produce negative humidity or negative rain, which are physical impossibilities. Clamping ensures the output remains within the $[0, 100]$ range for humidity. Furthermore, the 0.3mm rain threshold addresses the problem of meteorological noise. SLR may project negligible precipitation values based on minor downward moisture trends; applying this floor filters false positives, ensuring the AI only predicts rain when the mathematical trend suggests measurable precipitation.

\subsection{Quantitative Accuracy and User Trust Metric}
Validation is performed via a customized linear penalty function:
\begin{equation}
A = \max(0, 100 - |T_{pred} - T_{actual}| \times 10)
\end{equation}
While scientific models often use Mean Absolute Error, we employ a ``Trust Score.'' A multiplier of 10 means every $1^\circ\text{C}$ of error reduces the AI's score by 10\%. This provides a clear benchmark: a $2^\circ\text{C}$ error is 80\% Accuracy (Good), whereas a $5^\circ\text{C}$ error results in 50\% Accuracy, signaling a failure of the linear trend to capture that day's atmospheric shifts.

\section{Enterprise Potential and Worldwide Context}
This project holds significant enterprise potential due to its horizontal scalability and low operational cost. Unlike traditional forecasting services that charge high fees for API calls to proprietary NWP models, this framework can be deployed as a microservice in private cloud environments. 

On a worldwide scale, tools like \textbf{Google GraphCast} and \textbf{IBM's The Weather Company} are pushing the boundaries of AI-driven meteorology. However, these tools are often inaccessible for smaller enterprises or developing regions due to their high resource requirements. Our framework aligns with the global movement toward \textbf{Explainable AI (XAI)} and \textbf{Edge Computing}, providing a viable alternative for users who require transparency and low-bandwidth solutions in areas where global model resolution remains insufficient.

\section{Future Applications}
The lightweight and modular nature of this framework allows for a wide array of future applications across various sectors:

\begin{enumerate}
    \item \textbf{Precision Agriculture:} Small-scale farmers can utilize this tool to predict localized micro-climate shifts. By understanding the linear trend of soil-level humidity, irrigation systems can be optimized to conserve water.
    \item \textbf{Smart Home Optimization:} The engine can be integrated into IoT-enabled HVAC systems. By predicting the 24-hour temperature trend, a smart thermostat can pre-cool or pre-heat a building more efficiently than reactive sensors.
    \item \textbf{Renewable Energy Forecasting:} The model can be adapted to predict solar irradiance and wind persistence, allowing micro-grid managers to handle battery storage before a projected decline in clear weather.
\end{enumerate}

\section{Conclusion}
The Weather AI Analytics engine serves as a foundation for hybrid modeling. By combining deterministic momentum with adaptive weighting and a transparent verification loop, the system offers a verifiable alternative to traditional forecasting. Future iterations will incorporate barometric pressure trends to improve the detection of non-linear weather shifts.

\begin{thebibliography}{9}
\bibitem{openmeteo}
Zippenfenig, P. (2023). \textit{Open-Meteo: High-resolution weather API for developers}. Available at: \url{https://open-meteo.com/}.

\bibitem{regression}
Montgomery, D. C., Peck, E. A., \& Vining, G. G. (2021). \textit{Introduction to Linear Regression Analysis}. Wiley Series in Probability and Statistics.

\bibitem{ewma}
Hunter, J. S. (1986). ``The Exponentially Weighted Moving Average.'' \textit{Journal of Quality Technology}, 18(4), 203-210.
\end{thebibliography}

\end{document}